\subsection{Análisis crítico y evaluación de riesgos}

La evaluación de la Tabla~\ref{tab:gestion_riesgos} indica que la infraestructura implementada reduce la carga operativa asociada a la administración de identidad, base de datos y archivos. Supabase Auth centraliza las sesiones, PostgreSQL mantiene los datos estructurados y Storage conserva los audios en repositorios privados. FastAPI aplica las reglas de negocio y autorización antes de acceder a estos servicios.

El riesgo de transcripción incorrecta está parcialmente controlado. La ruta local permite editar el resultado antes de enviarlo y la ruta remota permite reproducir y confirmar la transcripción. Sin embargo, el ASR todavía requiere mediciones con vocabulario clínico, ruido ambiental y dispositivos representativos.

Los riesgos vinculados con NER, clasificación y generación de borradores permanecen asociados a la integración funcional de estos módulos. Aunque la arquitectura separa privacidad, procesamiento y validación, Hito 4 debe demostrar que el texto se sanitiza antes del LLM y que ningún contenido se publica sin HITL.

La integración hospitalaria simulada permite evaluar los contratos con datos sintéticos, pero todavía debe demostrar validación, normalización, idempotencia y manejo de errores mediante pruebas reproducibles. La sincronización offline y la resolución de conflictos continúan pendientes, por lo que una pérdida de conectividad puede interrumpir el registro. También debe comprobarse la cobertura de auditoría sobre accesos, descargas y cambios de permisos. Finalmente, la adopción no ha sido validada con usuarios reales; el protocolo definido deberá aplicarse antes de concluir que la interfaz se ajusta al trabajo clínico.
